% !TEX root = ../main.tex
%

\section{Conclusion}
While recent studies predominantly rely on proprietary LLMs for simulations, we argue that this choice is both excessive in terms of the required capabilities and unscalable due to high inference costs. To support this position, we demonstrated that experiments using small open-source models can reduce inference costs by more than forty-four, while still observing important, emergent social properties and producing relatively convincing discussions. Downscaling the models required the development of a unified, task-agnostic theoretical framework for \ac{sdg}, in which we identified key components and requirements, and proposed an evaluation schema to validate these elements.

Despite ongoing skepticism regarding the use and generalizability of LLM simulations, we showed that \ac{sdg} systems can yield meaningful insights by uncovering a critical limitation of LLM facilitators: their tendency to constantly intervene in discussions. Furthermore, we demonstrated that different facilitation strategies significantly influence the patterns of facilitation observed. Overall, our work represents a methodological shift in the design and execution of \ac{sdg} experiments. Our contributions include:
(1) a theoretical framework for designing, documenting, and evaluating \ac{sdg} systems,
(2) an open-source Python framework implementing this methodology,
(3) a cost-comparison methodology for social science experiments,
(4) an exploration of available models and algorithms,
(5) inherent limitations and a promising direction in LLM facilitation, and
(6) a large-scale dataset of synthetic discussions usable for \ac{sdg} analysis and LLM facilitator finetuning. 


\section{Discussion}

\subsection{Limitations}
Our research is fundamentally limited by the lack of robust evaluation metrics for synthetic text. While the diversity metric used throughout this study represents a meaningful step forward, we acknowledge that content variety is only a heuristic for measuring similarity to human discussion. More broadly, the field of \ac{gabm} still lacks reliable evaluation strategies that do not rely on extensive, specialized, and costly human annotation.

Our simulations are further constrained by the stereotypical and repetitive patterns of speech that many instruction-tuned models exhibit. The scale of our ablation experiments also prevents us from benchmarking proprietary models (\S\ref{sec:scalability:cost}; Table~\ref{tab:cost}), although we have shown that small models are sufficient for our simulations. As such, our findings remain unchanged. Finally, while our study focuses on monetary cost as the primary constraint faced by researchers, other factors such as computational speed and ease of use should also be considered in future work. 


\subsection{Future work}
Future work should focus on developing more robust methods for evaluating synthetic discussions and replicating social behaviors using LLMs. This includes exploring additional ways to quantify the similarity between human and LLM-generated discussions, as well as establishing standardized benchmarks to assess the utility of different models in the context of \ac{sdg} systems.

We believe that \ac{sdg} represents a promising research direction with broad applicability beyond LLM facilitation. For instance, it could be used to study the impact of state-actors leveraging LLM agents to influence public narratives, spread misinformation, or promote propaganda. Qualitative extensions of this work could involve the use of a dimorphic setup to distinguish between participants and facilitators, as discussed in \S\ref{sec:scalability:experimental}. Additionally, employing non-finetuned models could lead to more ""authentic"" generation of dialogue, as outlined in \S\ref{sec:scalability:model}. Furthermore, the development of sophisticated turn-taking algorithms that grant agency to LLM participants could enhance the realism and interactivity of these synthetic systems, as described in \S\ref{sec:methodology:management:turn-taking}.

Finally, while our research highlights that changing facilitation strategies significantly influences the behavior of LLM facilitators, the specific nature of these changes should be further explored through replication studies involving human participants. 


\subsection{Ethical considerations}
Synthetic discussions involving LLMs could be exploited by malicious actors to train them at performing unethical tasks \cite{majumdar_2024_nefarious,MARULLI20245340,li_2025_vulnerable}. Additionally, our research includes elements of LLM jail-breaking using conversational context, which could be used to poison discussions using safety-aligned models, although ongoing research is addressing these vulnerabilities \cite{wang_2025_risk}. Furthermore,the use of LLMs inherently risks skewing moderation systems towards the predominant demographics best represented in their training data. \ac{sdb} prompts are a necessary but insufficient step towards avoiding this \cite{rossi_2024,anthis_2025,burton2024large}. 

\subsection{AI use statement}
We have used small, open source LLMs (primarily LLaMa-8b and Qwen-7B) to help in small style changes throughout this document. We also used them in conjunction with a GPT-5 model \cite{openai2024gpt4technicalreport} to write initial documentation for our framework, as well as to generate some of the code for the plots presented in this paper. All such outputs have been checked by the authors.